\documentclass[11pt]{article}
\usepackage[margin=1in]{geometry}
\usepackage[T1]{fontenc}
\usepackage[utf8]{inputenc}
\usepackage{amsmath,amssymb,amsthm}
\usepackage{xcolor}
\usepackage{microtype}
\usepackage[colorlinks=true,linkcolor=blue!55!black,citecolor=blue!55!black,urlcolor=blue!55!black]{hyperref}
\newcommand{\B}{\mathbb B}
\newcommand{\C}{\mathbb C}

\title{Literature status: the tight-fitting conjecture in arXiv:2502.08406}
\author{Answer completed by arXiv:2607.20055v2 and its cited sharp-contraction results}
\date{11 August 2026}

\begin{document}
\maketitle

\begin{center}
\fbox{\parbox{0.92\textwidth}{\raggedright\textbf{Status: fully answered in later literature.}
Conjecture~22 of Bao--Ma--Yan--Zhu is now a consequence of Kalaj's proof of the
Gromov--Ros conjecture, arXiv:2607.20055v2, together with the source paper's own
embedding and compactness theorems.  Kalaj's Corollary~7.2 gives the sharp
Bergman contraction and all equality cases; Corollary~7.3 gives the Hardy
endpoint.  A one-line intermediate-norm argument below also gives all Hardy
endpoint equality cases.}}
\end{center}

\section{Question and normalization dictionary}

Conjecture~22 of arXiv:2502.08406 asks whether, for $a,b>-1$, the only tight
fittings on $\B^n\subset\C^n$ are
\begin{align*}
 A_a^p\subset A_b^q &: \quad p<q,\qquad (n+1+a)q=(n+1+b)p, \\
 H^p\subset A_a^q &: \quad p<q,\qquad nq=(n+1+a)p,
\end{align*}
and whether equality of the two norms occurs exactly for
\begin{align*}
 f(z)&=\frac{c}{(1-\langle z,w\rangle)^{2(n+1+a)/p}}
 &&\text{in the first line},\\
 f(z)&=\frac{c}{(1-\langle z,w\rangle)^{2n/p}}
 &&\text{in the second line},
\end{align*}
where $c\in\C$ and $w\in\B^n$.

Kalaj uses invariant measure
\[
 d\mu_n(z)=\frac{dv(z)}{(1-|z|^2)^{n+1}}
\]
and writes $\mathcal A_A^p$ for the normalized norm with integrand
$|f|^p(1-|z|^2)^A\,d\mu_n$, where $A>n$.  Thus
\[
 A=n+1+a,\qquad B=n+1+b,\qquad
 \mathcal A_A^p=A_a^p,\qquad \mathcal A_B^q=A_b^q.
\]
In this notation the first critical line is $q/p=B/A$ and the second is
$q/p=A/n$.

\section{Later theorem and the two contractions}

Li and Su~\cite{LiSu} had proved the required higher-dimensional convex
contraction conditionally on geodesic balls being all isoperimetric regions in
complex hyperbolic space.  Kalaj~\cite{Kalaj} proves exactly this missing input
in Theorem~6.9 for every $n\ge2$.  His Theorem~7.1 therefore makes the sharp
convex Bergman contraction unconditional.  In particular, Corollary~7.2 says
that for $A,B>n$ and $q/p=B/A>1$,
\[
 \|f\|_{\mathcal A_B^q}\le \|f\|_{\mathcal A_A^p},
\]
with equality, for nonzero $f$, exactly for
\[
 f(z)=c(1-\langle z,w\rangle)^{-2A/p}.
\]
The one-dimensional inequality is Kulikov's theorem; the strict-convexity
equality principle cited by Kalaj (Nicola--Riccardi--Tilli, Theorem~5.2) gives
the same equality family.  This proves the analytic assertion in part~(a) of
Conjecture~22 in every dimension.

For the Hardy endpoint, put $r=p/n$ and $A=n+1+a$.  The critical identity
$nq=Ap$ is exactly $q=rA$.  Kalaj's Corollary~7.3, the now-unconditional Hardy
endpoint of the Li--Su chain, gives
\[
 \|f\|_{\mathcal A_A^{rA}}\le \|f\|_{H^{rn}},
\]
which is precisely $\|f\|_{A_a^q}\le\|f\|_{H^p}$.  The corresponding disc
statement is again contained in Kulikov's one-dimensional result.

\section{Endpoint equality cases}

The full Hardy equality classification follows without any additional
geometric argument.  Choose an arbitrary $C$ with $n<C<A$.  The contractive
chain gives
\[
 \|f\|_{\mathcal A_A^{rA}}
 \le \|f\|_{\mathcal A_C^{rC}}
 \le \|f\|_{H^{rn}}.
\]
If the two endpoint norms are equal, then equality holds in the first
inequality.  The Bergman equality statement just quoted, applied with source
parameters $(C,rC)$ and target parameters $(A,rA)$, forces
\[
 f(z)=c(1-\langle z,w\rangle)^{-2C/(rC)}
     =c(1-\langle z,w\rangle)^{-2/r}
     =c(1-\langle z,w\rangle)^{-2n/p}.
\]
Conversely these kernels are extremal, as already observed in the source paper
(or by ball-automorphism invariance).  The case $f=0$ is included by taking
$c=0$.  This proves exactly the extremal assertion in part~(b).

\section{Why this proves ``if and only if'' for tight fittings}

It remains only to combine the sharp inequalities with results already proved
in arXiv:2502.08406.  Its Theorems~A--C show that a noncompact proper embedding
with $a,b>-1$ can occur in the two displayed families only on the stated
critical line and with $p<q$.  Conversely, on that critical line the embedding
is proper and noncompact.  The later sharp inequalities make it contractive;
constants show that its operator norm is exactly one.  Hence it is a tight
fitting.  This proves both directions and both equality classifications of
Conjecture~22.

\section{Verification note}

The source question was checked in Conjecture~22 of arXiv:2502.08406.  The
later status was checked against arXiv:2607.20055v2 (revised 7 August 2026),
especially Theorem~6.9, Theorem~7.1, Corollary~7.2, and Corollary~7.3, and
against the conditional theorem and parameter convention in
arXiv:2411.01911.  The parameter shift $A=n+1+a$ is forced by changing from
Euclidean weighted volume to invariant volume, and makes both critical lines
identical term by term.

\begin{thebibliography}{4}
\bibitem{Source}
G.~Bao, P.~Ma, F.~Yan, and K.~Zhu,
\emph{Embedding and compact embedding between Bergman and Hardy spaces},
arXiv:2502.08406; J. Geom. Anal. 36 (2026).

\bibitem{LiSu}
X.~Li and G.~Su,
\emph{Contraction property on complex hyperbolic ball},
arXiv:2411.01911, 2024/2025.

\bibitem{Kalaj}
D.~Kalaj,
\emph{The Gromov--Ros conjecture in complex hyperbolic space},
arXiv:2607.20055v2, 2026.

\bibitem{NRT}
F.~Nicola, F.~Riccardi, and P.~Tilli,
\emph{The Wehrl-type entropy conjecture for symmetric $SU(N)$ coherent states:
cases of equality and stability}, arXiv:2412.10940, especially Section~5.
\end{thebibliography}

\end{document}
